% ==========================================
% SECTION 3: CONSTANT PROPAGATION
% ==========================================
\section{Constant Propagation}

Constant Propagation is a global dataflow analysis designed to discover variables that are guaranteed to hold a single
constant value at specific program points across all possible execution paths. By tracking pairs of $<$variable, constant$>$, 
this analysis enables the compiler to substitute variable uses with their actual values and statically fold 
binary expressions whose operands are known constants, thereby optimizing execution efficiency.

\begin{table}[H]
\centering
\renewcommand{\arraystretch}{1.5}
\begin{tabular}{ll}
\toprule
\textbf{Parameter} & \textbf{Definition} \\
\midrule
\textbf{Domain} & Sets of pairs $<variable, value>$\\
\textbf{Direction} & Forward \\ & $in[b] = \wedge(out[pred(b)])$ \\ & $out[b] = f_B(in(b)]$ \\
\textbf{Transfer Function} & $f_b(x) = Gen_b \cup (x \setminus Kill_b)$ \\
%in[b] = Gen_b \cup (out[b] \setminus Kill_b)$ \\
\textbf{Meet Operation ($\wedge$)} & Intersection ($\cap$) \\
\textbf{Boundary Condition} & $in[entry] = \emptyset$, $out[entry] = \emptyset$ \\
\textbf{Initial Interior Points} & $in[b] = U$, $out[b] = \emptyset$ \\
\bottomrule
\end{tabular}
\caption{Contant Propagation parameters}
\end{table}

\begin{algorithm}[h!]
\caption{Constant Propagation (TBD)}
\begin{algorithmic}[1]
\State \Comment{Algoritmo da definire...}
\end{algorithmic}
\end{algorithm}